In principle, each page of a paper is numbered. The numbering begins with the
cover page, in the case of book publications with the first page of the title.
The book cover is not numbered. On certain pages, e.g. the cover page or on
cover page or on blank pages, no page number is printed. For more extensive
works, the numbering starts again at~1 from the first chapter or the preface.
The preceding pages (table of contents, etc.) are then numbered for the sake of
distinction, are then numbered with Roman numerals for the sake of distinction.
The numbering is switched with the \LaTeX command \verb+\pagenumbering+, as
shown in Listing~\ref{lst:final:pagenumber}.
The page counter is automatically reset to~1.

\lstinputlisting[language={[LaTeX]{TeX}},style=block,float,caption={switch page numbering},label={lst:final:pagenumber}]{code/final_pagenumber.en.tex}

In the document class~\texttt{book} the change of the numbering is done
automatically by the commands \verb+\frontmatter+, \verb+\mainmatter+ and
\verb+\backmatter+.

A separate numbering scheme for attachments is normally not customary.
If annexes are to be numbered in the form A-1, A-2, \dots, B-1, B-2, \dots\ must
be numbered according to Listing~\ref{lst:final:pagenumber_appendix}.
The construct of \verb+\pagenumbering+ and \verb+\renewcommand+ must be repeated
with each appendix. Corresponding adjustments must then also be made for the
the following directories, such as the bibliography, glossary, etc.
Continuous Arabic numbering is the more elegant solution here.

\lstinputlisting[language={[LaTeX]{TeX}},style=block,float,caption={Separate page numbering for appendices},label={lst:final:pagenumber_appendix},morekeywords={chapter}]{code/final_pagenumber_appendix.en.tex}
